\section{Discussion}
\label{sec:discussion}

The predicate-level intervention results of Section~\ref{sec:experiments} support the central claim of this paper: not all predicate extraction errors affect the final legal decision equally, and their impact depends systematically on where the predicate sits in the symbolic rule structure. This section interprets what that divergence means for evaluating and deploying neuro-symbolic legal reasoning pipelines, discusses model-specific reliability differences it exposes, and states the limitations of the present study.

\subsection{Extraction Difficulty Is Not Decision Importance}
\label{subsec:interpretation}
Standard practice evaluates an LLM predicate extractor by its F1-score against a gold predicate table, implicitly treating extraction difficulty as a proxy for how much a given predicate's errors matter\cite{tan2026diagnosing}. Table~\ref{tab:dir} shows this proxy fails in both directions. \textit{is\_written\_assertion} is simultaneously the hardest predicate to extract (F1 between 0.286 and 0.429) and the one whose errors are most likely to flip the final decision (DIR between 0.700 and 0.900): here, difficulty and importance coincide, and this predicate deserves priority for extractor improvement. \textit{is\_non\_verbal\_assertive\_conduct} has comparably low F1 (0.593 - 0.667) but near-zero DIR (0.000 - 0.091): its errors are difficult to extract but structurally redundant, because the disjunctive assertion clause of Equation~\ref{eq:hearsay-rule} is almost always already satisfied by \textit{is\_oral\_assertion} or \textit{is\_written\_assertion} on the same instance. Conversely, \textit{is\_out\_of\_court} and \textit{is\_introduced\_to\_prove\_truth} - the two conjunctive gates in the rule - are extracted comparatively well, yet their rare errors carry disproportionate weight (DIR up to 1.000 and 0.875, respectively), because a conjunctive condition has no disjunctive backup to absorb a mistake. A practitioner who allocates extractor-improvement effort purely by predicate F1 would over-invest in \textit{is\_non\_verbal\_assertive\_conduct} and under-invest in the two conjunctive gates - exactly the misallocation that rule-aware error attribution is designed to prevent.

\subsection{From Diagnosis to Verification: Rule-Aware Triage}
\label{subsec:triage}
This divergence has a direct practical use that connects back to the verification role PYTHEN is intended to play in the broader neuro-symbolic pipeline: DIR can prioritize which extracted predicates are worth verifying, either by a human reviewer or by a secondary model call, before a decision is trusted. Rather than auditing all five predicates uniformly, or auditing whichever predicates the extractor is least confident about, a rule-aware verification policy would flag only extracted predicates that are both erroneous-prone and high-DIR for that instance's rule path - concretely, \textit{is\_written\_assertion} and \textit{is\_introduced\_to\_prove\_truth} in the hearsay setting, and, whenever the assertion clause is otherwise unsatisfied, \textit{is\_out\_of\_court}. Because low-DIR predicates such as \textit{is\_non\_verbal\_assertive\_conduct} rarely change the outcome even when wrong, they can be deprioritized without materially increasing the risk of an incorrect final decision. This reframes predicate-level verification from an undifferentiated checklist into a targeted, rule-path-dependent audit, reducing verification cost without a proportional loss of reliability - one concrete way the symbolic layer can support explanation and verification of LLM output in high-stakes legal settings.

\subsection{Model-Specific Reliability Profiles}
\label{subsec:model-differences}
Table~\ref{tab:taxonomy} shows that the same nominal error type carries different practical risk across models. DeepSeek V3.2 and GPT-OSS 120B's hallucinated predicates (E1) are decision-critical only 59.5\% and 46.9\% of the time, respectively - meaning a large share of their false positives are structurally absorbed and never reach the output. Gemma 4 31B breaks this pattern: despite hallucinating about as often in absolute terms, 96.2\% of its E1 errors are decision-critical. This is not visible from aggregate accuracy or F1 alone, where Gemma 4 31B is simply the strongest model on every metric (Table~\ref{tab:predicate-extraction}); the taxonomy reveals that when Gemma 4 31B does err, it tends to err on predicates the decision actually depends on, rather than on redundant ones. This suggests that predicate-level F1 rankings and decision-level reliability rankings need not coincide, and that a model selection or deployment decision informed only by the former risks underestimating how consequential a given model's remaining errors are in practice.

\subsection{Limitations}
\label{subsec:limitations}
\textbf{Single task and rule structure.} All results in this paper come from the \texttt{hearsay} subset of LegalBench, whose rule tree is a single shallow disjunction gated by two conjunctive conditions (Equation~\ref{eq:hearsay-rule}). LegalBench and the broader legal NLP literature contain tasks with substantially different rule shapes - deeper hierarchies\cite{lan2026hscodecomp}, multiple sequential gates\cite{fan2026lexam}, exception and override clauses\cite{cooper2026defab}, or predicates with numeric thresholds\cite{sadowski2025verifiable} - that have not been evaluated here. The specific divergence pattern we observe (a hard-to-extract predicate that is also high-impact, alongside a hard-to-extract predicate that is low-impact) is a property of this particular rule structure, and it remains an open question whether rule-aware error attribution surfaces the same kind of pattern, a different one, or none at all, on rule trees of a different shape. We treat this paper's findings as a case study establishing the methodology and its value on one task, not as a generalized claim about neuro-symbolic legal pipelines at large.

\textbf{Small-sample DIR estimates.} DIR is computed only over the errors a model actually makes, and for some predicate-model pairs this is a small count (e.g., 3 errors for \textit{is\_out\_of\_court} under Gemma 4 31B, 5 for \textit{is\_oral\_assertion}). A DIR of 1.000 estimated from 3 errors is directionally informative but statistically fragile, and should be read alongside the error count reported in Table~\ref{tab:dir} rather than as a precise probability.

\textbf{Single-predicate intervention does not capture interaction effects.} By construction (Equation~\ref{eq:intervention}), our intervention corrects one erroneous predicate at a time while holding all others - including other errors - fixed. When an instance has two or more simultaneous predicate errors, individually correcting either one may fail to flip the decision even though jointly correcting both would; such instances are undercounted as "no critical predicate found" under our procedure, even though the errors are jointly responsible for the wrong decision. A full combinatorial intervention over all subsets of $E_i$ would resolve this at the cost of exponential evaluation cost in $|E_i|$, which we did not pursue given that multi-error instances are a minority in our data.

\textbf{Gold predicate annotation and rule soundness are validated only on this dataset.} The oracle experiment (Section~\ref{subsec:oracle}) shows that the PYTHEN rule tree reproduces every gold label when given correct predicates for the \texttt{hearsay} subset, and the gold predicate table itself was manually annotated by the authors. Both the rule tree's logical soundness and the predicate annotations are therefore validated only within the scope of this dataset, not against independent legal experts or held-out cases.

\subsection{Future Work}
\label{subsec:future-work}
The most direct extension is to repeat the predicate-level intervention procedure on additional LegalBench tasks with structurally different rule trees - particularly ones involving exception clauses or sequential gating - to test whether the extraction-difficulty/decision-impact divergence observed here is a general property of neuro-symbolic legal pipelines or specific to the hearsay rule's shape. A second direction is to use DIR directly as a training or data-collection signal: weighting fine-tuning examples, few-shot demonstrations, or active-learning acquisition by a predicate's DIR rather than its base rate or current F1 would target extractor improvement at the errors most likely to change outcomes. Third, the rule-aware triage strategy outlined in Section~\ref{subsec:triage} could be implemented and evaluated as an actual human-in-the-loop verification interface, measuring whether DIR-prioritized review recovers most of the reliability of full predicate verification at a fraction of the review cost. Finally, extending the intervention procedure from single-predicate to joint (multi-predicate) corrections would address the interaction-effect limitation above and provide a more complete account of which \emph{combinations} of errors are jointly responsible for incorrect legal decisions.
